\chapter{Nachweis der Radioaktivität in einer Probe}

Im abschließenden Versuchsteil wurde mit dem HPGe-Detektor das Gammaspektrum einer Bodenprobe aufgenommen. 
Durch den Vergleich der gemessenen Energien charakteristischer Linien mit bekannten Zerfallsdaten radioaktiver Nuklide können Rückschlüsse auf die in der Probe enthaltenen radioaktiven Stoffe gezogen werden.

\section{Durchführung}

Die untersuchte Probe bestand aus etwa einem Kilogramm handelsüblicher Blumenerde auf Basis von Hochmoortorf. 
Das Material befand sich in einem dünnwandigen Kunststoffbeutel und wurde unmittelbar vor dem HPGe-Detektor positioniert, um die Nachweiswahrscheinlichkeit der emittierten $\gamma$-Quanten zu maximieren.

Da für natürliche Bodenproben lediglich geringe Aktivitäten zu erwarten sind, wurde eine Messdauer von 15,h,45,min gewählt. 
Durch die lange Messzeit wird die statistische Unsicherheit der registrierten Ereigniszahlen reduziert. 
Für die relative statistische Unsicherheit gilt $\frac{\Delta N}{N}=\frac{1}{\sqrt{N}}$,
sodass größere Ereigniszahlen zu einer verbesserten Genauigkeit der späteren Auswertung führen.

Zur Verringerung des Einflusses äußerer Strahlungsquellen wurden Probe und Detektor durch Bleiblöcke abgeschirmt (vgl. Abb.\ref{fig:aufbau_1_2}). 
Zusätzlich war vorgesehen, ein separates Hintergrundspektrum unter identischen Messbedingungen aufzunehmen, um die Umgebungsstrahlung vom Probenspektrum subtrahieren zu können. 
Hierfür sollte die Probe aus dem Aufbau entfernt werden, während Geometrie und Messdauer unverändert bleiben.

Die zugehörigen Messdaten standen jedoch für die vorliegende Auswertung nicht zur Verfügung. 
Da es sich um eine Langzeitmessung handelte und der Versuchsaufbau bereits vor Bereitstellung der Daten abgebaut wurde, konnte keine Hintergrundkorrektur durchgeführt werden. 
Die nachfolgende Analyse basiert daher ausschließlich auf dem gemessenen Rohspektrum der Bodenprobe.


\section{Extraktion der $\gamma$-Energien}

Das aufgenommene Spektrum der Bodenprobe ist in Abb.~\ref{fig} dargestellt. 
Insgesamt konnten 16 deutlich ausgeprägte Linien identifiziert werden. 
Zur Bestimmung ihrer Lage wurde jede Struktur mit einer Gaußfunktion unter Berücksichtigung eines linearen Untergrunds angepasst. 
Die verwendete Fitfunktion wurde bereits in Gleichung~\ref{gl:} eingeführt.

Die angepassten Kurven sind gemeinsam mit dem gemessenen Spektrum in Abb.~\ref{fig:erde_spektrum} dargestellt. 
Aufgrund der geringen Breite einzelner Peaks wurden zusätzlich Detaildarstellungen der jeweiligen Fits angefertigt. 
Die ermittelten Schwerpunktlagen der Peaks sind zusammen mit den daraus berechneten Energien in Tabelle~\ref{tab:erde_peaks} aufgeführt.

Im niederenergetischen Bereich des Spektrums ist ein deutlicher Anstieg des Untergrunds zu erkennen, dessen Intensität mit zunehmender Energie kontinuierlich abnimmt. 
Dieser Effekt kann auf verbleibende Untergrundstrahlung sowie auf charakteristische Röntgenstrahlung der verwendeten Bleiabschirmung zurückgeführt werden.

Die Umrechnung der Peakpositionen in Energien erfolgte mithilfe der zuvor bestimmten Energiekalibration des HPGe-Detektors gemäß Gleichung~\ref{gl.}. Die hierfür verwendeten Kalibrationsparameter sind in Abschnitt~\ref{} angegeben.


\begin{figure}[H]
    \centering
    \includegraphics[width=\textwidth]{Ge_Erde_ramen_gauss.png}
    \caption{Gemessenes HPGe-Spektrum der untersuchten Blumenerdeprobe. 
    Die eingezeichneten Kurven entsprechen den Gaußanpassungen der identifizierten Peaks. 
    Eine Korrektur durch Subtraktion eines Hintergrundspektrums wurde nicht durchgeführt.}
    \label{fig:erde_spektrum}
\end{figure}


\begin{table}[H]
\centering
\caption{Identifizierte $\gamma$-Linien der Blumenerdeprobe ohne abzug der Untergrundes.}
\label{tab:erde_peaks}
\begin{tabular}{cccc}
\toprule
Peak & $E_\gamma$ / keV & Kanal $\mu$ & Zuordnung \\
\midrule
1  & $72.77 \pm 0.03$   & $753.67 \pm 0.28$   & Pb-K$\alpha$ \\
2  & $74.96 \pm 0.02$   & $776.37 \pm 0.16$   & Pb-K$\alpha$ \\
3  & $84.79 \pm 0.03$   & $878.35 \pm 0.33$   & Pb-K$\beta$ \\
4  & $87.36 \pm 0.06$   & $905.03 \pm 0.62$   & Pb-K$\beta$ \\
5  & $238.67 \pm 0.02$  & $2474.84 \pm 0.18$  & $^{212}$Pb \\
6  & $295.23 \pm 0.04$  & $3061.66 \pm 0.44$  & $^{214}$Pb \\
7  & $338.29 \pm 0.06$  & $3508.37 \pm 0.62$  & $^{228}$Ac \\
8  & $351.95 \pm 0.02$  & $3650.12 \pm 0.22$  & $^{214}$Pb \\
9  & $510.94 \pm 0.04$  & $5299.70 \pm 0.43$  & Annihilation \\
10 & $583.19 \pm 0.03$  & $6049.29 \pm 0.29$  & $^{208}$Tl \\
11 & $609.39 \pm 0.02$  & $6321.09 \pm 0.25$  & $^{214}$Bi \\
12 & $911.22 \pm 0.03$  & $9452.59 \pm 0.30$  & $^{228}$Ac \\
13 & $964.80 \pm 0.13$  & $10008.44 \pm 1.30$ & $^{228}$Ac \\
14 & $969.01 \pm 0.05$  & $10052.17 \pm 0.53$ & $^{228}$Ac \\
15 & $1120.41 \pm 0.06$ & $11622.98 \pm 0.58$ & $^{214}$Bi \\
16 & $1460.86 \pm 0.02$ & $15155.14 \pm 0.10$ & $^{40}$K \\
\bottomrule
\end{tabular}
\end{table}

\section{Zuordnung der $\gamma$-Energien zu ihren Quellen}

Die aus den Fitparametern bestimmten Energien wurden mit den im NuDat-3-Datensatz aufgeführten Zerfallsdaten verglichen, um die beobachteten Linien geeigneten Nukliden zuzuordnen \cite{NuDat3}. Die im Folgenden genannten Intensitäten stammen ebenfalls aus den dort aufgeführten Kernzerfallsdaten.

Die Peaks 1 bis 4 lassen sich mit hoher Wahrscheinlichkeit charakteristischen Röntgenlinien des Bleis zuordnen. Die entsprechenden Literaturwerte betragen 72,80542(24)\,keV, 74,97011(17)\,keV, 84,45045(60)\,keV und 84,93908(34)\,keV. 
Die gute Übereinstimmung mit den gemessenen Energien spricht dafür, dass diese Linien überwiegend durch die verwendete Bleiabschirmung verursacht werden \cite{NIST_Xray}.

Peak 5 wird dem Nuklid $^{212}$Pb aus der natürlichen Thorium-Zerfallsreihe zugeordnet. Die Literaturenergie der dominanten Linie beträgt 238,632(2)\,keV.
Aufgrund der vergleichsweise hohen relativen Intensität von 43,6(5)\,\% sowie der Halbwertszeit von etwa 10\,h ist diese Identifikation sehr plausibel \cite{Auranen2020_212Pb}.

Die Peaks 6 und 8 stimmen ausgezeichnet mit den bekannten $\gamma$-Linien von $^{214}$Pb überein. 
Die Literaturwerte liegen bei 295,224(2)\,keV beziehungsweise 351,9320(21)\,keV. 
Beide Übergänge besitzen hohe Emissionswahrscheinlichkeiten von 18,47(11)\,\% und 35,72(24)\,\%, wodurch ihre Beobachtung in natürlichen Proben zu erwarten ist \cite{Zhu2021_214Pb}.

Die Linie bei etwa 511\,keV (Peak 9) entspricht der charakteristischen Annihilationsstrahlung eines Elektron-Positron-Paares. 
Solche Photonen können als Folge verschiedener Zerfallsprozesse oder sekundärer Wechselwirkungen hochenergetischer $\gamma$-Quanten entstehen. 
Eine eindeutige Zuordnung zu einem einzelnen Radionuklid ist daher nicht möglich.

Die Peaks 7, 12, 13 und 14 werden dem Nuklid $^{228}$Ac aus der Thorium-Zerfallsreihe zugeschrieben. 
Die entsprechenden Literaturenergien betragen 338,320(3)\,keV, 911,204(4)\,keV, 964,766(10)\,keV und 968,971(17)\,keV, mit relativen Intensitäten 11,27(19) \%, 25,8(4) \%, 4,99(9) \% und 15,8(3) \%.. 
Insbesondere die Linien bei 338\,keV, 911\,keV und 969\,keV gelten als charakteristische Signaturen dieses Nuklids \cite{Abusaleem2014_228Ac}.

Peak 10 bei etwa 583\,keV wird dem Nuklid $^{208}$Tl aus der Thorium-Zerfallsreihe zugeordnet.
Die Literaturenergie beträgt 583,1118(20)\,keV oder 583,1855(20)\,keV. Mit relativen Intensitäten von 0,00415(6)\,\% beziehungsweise 0,00344(5)\,\% ist es warscheinlich das es eine überlappung ist mit gewichtung auf das erste und stimmt aber sehr gut mit dem gemessenen Wert überein. 
Die Linie zählt zu den stärksten Übergängen dieses Nuklids und ist daher in natürlichen Proben häufig nachweisbar \cite{Martin2007_208Tl}.

Die Peaks 11 und 15 lassen sich eindeutig dem Nuklid $^{214}$Bi zuordnen aus der Uran-Radium-Reihe (Halbwertszeit 19,71(2) min). 
Die entsprechenden Literaturwerte betragen 609,321(7)\,keV und 1120,294(6)\,keV. 
Mit relativen Intensitäten von 45,44\,\% beziehungsweise 14,90(8)\,\% gehören diese Übergänge zu den dominierenden Linien des Nuklids \cite{Zhu2021_214Pb}.

Peak 16 wird dem natürlich vorkommenden Isotop $^{40}$K zugeordnet. 
Die charakteristische $\gamma$-Linie besitzt eine Literaturenergie von 1460,820(5)\,keV und eine relative Intensität von 10,66(17)\,\% \cite{Chen2017_40K}. 
Da $^{40}$K keine weiteren ausgeprägten $\gamma$-Linien emittiert, ist diese Identifikation eindeutig.
Bereits in früheren Hintergrundmessungen ohne zusätzliche Probe konnte im Bereich dieser Energie eine schwache Struktur beobachtet werden (vgl,Abb.\ref{fig:Ge_Cobalt}). 
Dies deutet darauf hin, dass ein Teil des Signals auf natürlich vorhandenes Kalium in der Umgebung oder auf Restaktivität des Messaufbaus zurückzuführen sein könnte. 
Aufgrund der fehlenden Langzeit-Hintergrundmessung lässt sich dieser Beitrag jedoch nicht quantitativ bestimmen.



Insgesamt zeigt die Untersuchung der Blumenerdeprobe das Vorhandensein mehrerer Nuklide aus den natürlichen Uran-Radium- und Thorium-Zerfallsreihen sowie des natürlich vorkommenden Isotops $^{40}$K. 
Erwartet worden war primär der Nachweis von Kalium, da dieses explizit als Bestandteil der verwendeten Düngemittel angegeben wird. 
Die zusätzlich beobachteten Zerfallsprodukte können plausibel durch weitere Bestandteile der Blumenerde erklärt werden, insbesondere durch Hochmoortorf sowie organische Materialien aus Garten-, Landschafts- und Forstwirtschaft.

Eine eindeutige Aussage über die genaue Herkunft der nachgewiesenen Nuklide ist anhand der vorliegenden Daten jedoch nicht möglich.

Grundsätzlich wäre auf Basis der bestimmten Peakflächen und unter Berücksichtigung der zuvor ermittelten energieabhängigen Nachweiseffizienz eine quantitative Aktivitätsbestimmung der einzelnen Nuklide möglich. 
Eine derartige Analyse würde jedoch den Umfang der vorliegenden Auswertung deutlich überschreiten.




